% Generated draft: requires mathematical and source review.
\documentclass[11pt]{article}
\usepackage[margin=25mm]{geometry}
\usepackage{fontspec}
\setmainfont{Latin Modern Roman}
\usepackage{amsmath,amssymb,mathtools,unicode-math}
\setmathfont{Latin Modern Math}
\usepackage{enumitem,needspace,xurl,hyperref}
\setlength{\parindent}{0pt}
\setlength{\parskip}{5pt}
\setlength{\emergencystretch}{3em}
\newcommand{\N}{\mathbb{N}}
\newcommand{\Z}{\mathbb{Z}}
\newcommand{\Q}{\mathbb{Q}}
\newcommand{\R}{\mathbb{R}}
\newcommand{\C}{\mathbb{C}}
\newcommand{\F}{\mathbb{F}}
\newcommand{\A}{\mathbb{A}}
\newcommand{\PP}{\mathbb{P}}
\newcommand{\E}{\mathbb{E}}
\newcommand{\eps}{\varepsilon}
\newcommand{\dd}{\,\mathrm{d}}
\begin{document}
\section*{UnsolvedMath ID: 8800008}
\textbf{Source problem code:} catalogue entry
\subsection{Definitions and mathematical statement}
The source asks to describe a cocycle $c_u$ explicitly, or equivalently to compute rapidly the invariants of certain local algebras split by generalized-dihedral extensions specified in the slides. The defining parameter $u$, cocycle construction and extensions are not reproduced.

In the standard local Brauer-group framework, let $K$ be a nonarchimedean local field and $L/K$ a finite Galois splitting extension with group $G$. A normalized multiplicative $2$-cocycle is a map $c:G\times G\to L^\times$ satisfying
\[
c(g,h)c(gh,k)=g(c(h,k))c(g,hk),\qquad c(1,g)=c(g,1)=1.
\]
Its crossed product has basis symbols $U_g$ over $L$, with $U_ga=g(a)U_g$ and $U_gU_h=c(g,h)U_{gh}$, and represents a central-simple-algebra Brauer class split by $L$. The local invariant in this convention lies in $\Q/\Z$ under $\operatorname{inv}_K:\operatorname{Br}(K)\to\Q/\Z$.

A finite generalized-dihedral group ordinarily has the form $A\rtimes\langle\tau\rangle$, where $A$ is abelian, $\tau^2=1$, and $\tau a\tau^{-1}=a^{-1}$. This identifies standard terminology only; the source-specific local fields, Galois realization, action and cocycle are absent. A splitting extension alone does not identify a unique Brauer class, so those data cannot be replaced by only the group name.

The requested output is an explicit formula or effective invariant computation for the specified $c_u$/algebra family. ``Fast'' needs an input representation, size measure and complexity target. The claimed equivalence of tasks also needs the actual algebra/cocycle presentation and reductions. The record supplies neither, so the general displayed framework is not a reconstruction of $c_u$ or an algorithm for it.
\subsection{Short English statement}
Compute the slide-specific generalized-dihedral cocycle and associated local algebra invariants efficiently. The cocycle parameter, field extensions, algebra presentation and computational input are missing.
\subsection{Sources}
\begin{itemize}
\item Local UnsolvedMath registry, permanent ID 8800008, exact source code catalogue entry, ``Explicit cocycles and invariants for split local algebras''; attributed to Gerhard Frey, no proposed year supplied. Registry status open, classification OPEN-TRIAGE.
\item \emph{IPAM Workshop on Number Theory and Cryptography---Open Problems} (2006), attachment 6219, PDF pages 36 and 44, Problem 7, \url{https://www.ipam.ucla.edu/programs/workshops/workshop-i-number-theory-and-cryptography-open-problems/?tab=schedule}. Supplied pdf-and-legacy-ppt visual-audit extraction accessed July 29, 2026. The provenance records inspection of 17 PDFs and two legacy PPT decks with repeated builds and solved definitions excluded; no PDF/PPT was accessed here.
\item The supplied summary mentions classical cyclic-algebra/Hasse-invariant techniques but reports no located citable resolution for this specific slide family. Its suggestion that standard methods may solve the problem is speculative and is not treated as an algorithm or proof.
\item Provenance flags current primary-literature status and redistribution terms NEEDS\_REVIEW.
\end{itemize}
\end{document}
